% ===========================================================================
%  The Threshold Does Not Transfer — standalone manuscript
%  Assembled from the section drafts in this directory. Section order follows
%  README.md; the deploy box sits between the introduction and the results by
%  decision, not by accident.
% ===========================================================================
\documentclass[a4paper,fleqn]{cas-dc}

\usepackage[numbers]{natbib}
\usepackage{amsmath,amssymb}
\usepackage{booktabs}
\usepackage{graphicx}
\usepackage{multirow}
\usepackage{array}
\graphicspath{{../}{../figures/}}

\newcommand{\macf}{macro-F\textsubscript{1}}

% The CAS class prints an "ORCID(s):" line even when no ORCID is supplied.
\makeatletter
\AtBeginDocument{\@ifundefined{printorcid}{}{\renewcommand{\printorcid}{}}}
\makeatother

\begin{document}

\let\WriteBookmarks\relax
\def\floatpagepagefraction{1}
\def\textpagefraction{.001}

\shorttitle{The threshold does not transfer}
\shortauthors{M. Kbaier et al.}

\title[mode=title]{The Threshold Does Not Transfer: Open-Set Test-Time Adaptation
Across Buildings}

\begin{abstract}
Deploying an activity recogniser in a new building means facing two shifts at once: the
sensors are in different places, and the occupant may do things the model was never
trained on. Open-set and universal domain adaptation address the second, and test-time
adaptation the first, but the methods for both are evaluated offline, with repeated passes
over a target set that is also the evaluation set. We study the combined shift under
constraints a deployment can meet: a source model trained once and frozen, a single online
pass over the unlabelled target stream in recording order, no source data, and no target
labels anywhere, including for the threshold that decides what counts as unknown.

Measured against a model trained on the target building at matched sample size, the
cross-building penalty on known activities is $1.6$ macro-F\textsubscript{1} points
entering one apartment and $36.0$ entering the other, a wider asymmetry than the
$0.6$ and $21.3$ of the closed-set case \cite{kbaier_closedset}. Four test-time methods
follow it: none improves on a frozen source model where there is no shift to recover.
One appears to, and the appearance is instructive. PETAL records the best open-set
H-score in one direction while classifying known activities roughly half as well, because
degrading a model uniformly lowers its confidence everywhere, pushes more windows below a
fixed threshold, and raises the unknown-detection term that H is bottlenecked by. The
standard open-set score is therefore maximised in part by damage, and pairing it with
known-class accuracy does not remove the problem; it only makes it visible.

The rejection rule fails for a different reason. Maximum softmax, entropy, maximum logit
and free energy are all at or below chance as novelty detectors here, with bootstrap
intervals excluding $0.5$ for three of the four: the model is as confident on activities it
never saw as on those it was trained on. A Mahalanobis distance to source class prototypes
is the only score above chance, and at a matched rejection rate it catches $1.58$ and
$1.17$ times as many unknowns as rejecting at random --- real, and far smaller than the
H-score comparison suggests, because two scores held at the same nominal percentile reject
very different fractions of a target stream.

That last observation is the main result. A threshold set to reject $1\%$ of source
validation rejects between $7\%$ and $63\%$ of the target stream, depending only on which
score carries it. Improving detection does not fix it: a source-only weighting rule that
gains $0.064$ AUROC in the direction with headroom \emph{loses} $0.11$ H at the threshold a
deployment would use, while the variant with the best H in both directions has
significantly worse AUROC in one. Detection quality and deployed open-set performance come
apart, so cross-domain threshold calibration is the binding constraint rather than one
difficulty among several. We recommend reporting the rejection rate, which is observable
without labels, beside every open-set score.
\end{abstract}

% ---------------------------------------------------------------------------


\author[1]{Mariem Kbaier}
\ead{kbaiermeryem@gmail.com}

\author[1]{Ons Abderrahim}

\author[1]{Manar Amayri}

\author[1]{Nizar Bouguila}
\cormark[1]
\ead{nizar.bouguila@concordia.ca}

\affiliation[1]{organization={Concordia Institute for Information Systems
  Engineering, Concordia University},
  city={Montreal}, state={QC}, country={Canada}}

\cortext[1]{Corresponding author.}

\begin{keywords}
Test-time adaptation \sep Open-set domain adaptation \sep Human activity recognition
\sep Smart buildings \sep Novelty detection
\end{keywords}

\maketitle

% ============================================================================
\section{Introduction}
\label{sec:os-intro}
% ---------------------------------------------------------------- motivation & setup
\subsection{Motivation and protocol}
\label{sec:os-setup}

Measured against a model trained on the target building at matched sample size, the
cross-building penalty for activity recognition in this pair of apartments is $0.6$
macro-F1 points entering $\text{S}_2$ and $21.3$ entering $\text{S}_1$: in one direction
there is essentially no distribution shift for an adaptation method to recover, and in the
other there is a substantial one. That asymmetry, rather than any single method's score,
is what decided the outcome of our closed-set study, and the methods tracked it --- losing
ground wherever there was no headroom and gaining only where there was.

This section asks whether a second, simultaneous shift changes that. We withhold two
activities, \texttt{toileting} and \texttt{cleaning}, from source training and treat them
as target-private, so the model faces a change of building and a change of label space at
once. The retained classes are \texttt{meal\_prep}, \texttt{watching\_tv} and
\texttt{music}; \texttt{cleaning} is placed in the unknown set deliberately, since at $113$
target windows it is not learnable in this data and leaving it among the known classes
would confound open-set performance with a sample-size limit. The experiment is a
$2\times2$ grid, $\{\text{S}_1\!\to\!\text{S}_2, \text{S}_2\!\to\!\text{S}_1\}$ crossed
with $\{$closed-set, open-set$\}$, and because withholding classes changes the source
model, every cell is given its own penalty measurement before any method is run.

Throughout we hold to the Tier-1 definition of test-time adaptation used in the
closed-set study: the source model is trained once and then frozen, with no further
multi-epoch training on target data; adaptation is online, in a single pass over the
unlabelled target stream, in recording order, with each batch predicted before the model
is updated on it; source data is unavailable during adaptation; and no target labels are
used anywhere, including for selecting hyperparameters and for calibrating the
unknown-rejection threshold, which is fixed from source validation before the target
stream is read. Any quantity a method relies on --- scalers, prototypes, thresholds ---
must be estimable without target labels. Where we report a figure that violates this, as
with the labelled ceiling of Section~\ref{sec:os-depth} and the matched-rejection-rate
comparison of Section~\ref{sec:os-maha}, it is labelled a diagnostic and is used to bound
what a method could achieve rather than to claim what one did.

The constraint matters because the reference implementations in this area do not observe
it. Section~\ref{sec:os-lead} documents four places where the released code for LEAD
leaves the test-time setting, including fifty offline epochs over the target set and an
evaluation loader that wraps the same data the model adapted on. Results obtained that way
are not comparable with results obtained under the constraints above, which is why every
number in this section is reported against our own Tier-1 source-only baseline rather than
against published figures.

% ---------------------------------------------------------------- what survives
\subsection*{What to deploy}
\label{sec:os-deploy}

This section records, before the analysis that justifies it, the one configuration in this
study that we would put in a building.

\begin{center}\small
\fbox{\begin{minipage}{0.92\columnwidth}
\textbf{Frozen source model, no adaptation, Mahalanobis rejection.}\\[2pt]
Train once on the source building's known activities. Do not adapt at test time. Reject a
window as unknown when its squared Mahalanobis distance to the nearest source class
prototype exceeds the $95$th percentile of that distance on the source validation split,
with the prototypes and the threshold both fixed before the target stream is read.\\[4pt]
H-score $0.5176$ entering $\text{S}_2$, $0.4289$ entering $\text{S}_1$.
\end{minipage}}
\end{center}

Every component of this is Tier-1: no target labels, no target statistics, one pass, and no
quantity estimated from the stream it is scored on. It is also, deliberately, a system
with no adaptation in it. Sections~\ref{sec:os-methods} to~\ref{sec:os-depth} report what
each of the parts we removed was worth, and none of them was worth keeping: the four
adaptation methods cost between $12$ and $27$ H-score points, the confidence-based
rejection scores are at or below chance, and the learned embedding contributes nothing the
raw features do not.

We state the recommendation first because the rest of this section is a sequence of
negative results, and a reader is entitled to know what we think survives them.


% ============================================================================
\section{Results}
\label{sec:os-res}
% ---------------------------------------------------------------- penalty grid
\subsection{The cross-building penalty in four cells}
\label{sec:os-penalty}

Withholding activities from source training changes the source model, so each cell of the
grid needs its own baseline before any method is run. We reuse the measurement of
the closed-set study \cite{kbaier_closedset}: hold the evaluation set fixed, vary only where the training
labels came from, match training-set size by subsampling, and group the split by recording
day. In the open-set cells the penalty is computed on known-class windows only, since a
source model that never saw the withheld activities cannot be asked to classify them.

% source: src/exp24_openset_gap.py -> results/exp24_openset_gap.csv, 3 seeds x 5 folds
\begin{table}[t]
\centering
\caption{Cross-building penalty, known-class \macf, matched training-set size, identical
held-out target days. Closed-set rows reproduce \cite{kbaier_closedset} on an
independently written harness.}
\label{tab:ospenalty}
\small
\begin{tabular}{@{}llccc@{}}
\toprule
Setting & Direction & Own days & Other building & Penalty \\
\midrule
Closed & $\text{S}_1\!\to\!\text{S}_2$ & $0.3497\pm0.148$ & $0.3434\pm0.177$ & $0.6$ \\
Closed & $\text{S}_2\!\to\!\text{S}_1$ & $0.4897\pm0.058$ & $0.2767\pm0.091$ & $21.3$ \\
Open   & $\text{S}_1\!\to\!\text{S}_2$ & $0.8082\pm0.133$ & $0.7926\pm0.137$ & $1.6$ \\
Open   & $\text{S}_2\!\to\!\text{S}_1$ & $0.9153\pm0.091$ & $0.5557\pm0.181$ & $\mathbf{36.0}$ \\
\bottomrule
\end{tabular}
\end{table}

Two points about how these numbers should be read.

The closed-set rows reproduce $0.6$ and $21.3$ points exactly, on a harness written
separately for this section. That is the only validation available to us that the
measurement is implemented correctly rather than merely reported consistently.

The small penalties are not small effects; they are absent ones. At $0.6$ and $1.6$ points
against between-fold standard deviations of $0.13$ to $0.18$, neither is distinguishable
from zero, and we read both as no measurable cross-building shift rather than as a shift
that happens to be slight. The large ones are the opposite: $21.3$ and $36.0$ points
against the same spread.

Absolute \macf{} is not comparable between the closed and open rows, since the open-set
task has three classes and the closed-set task five, and the three retained activities are
the better-populated ones. Only the penalties are comparable, and the comparison is that
restricting the label space \emph{widens} the asymmetry rather than evening it out:
$0.6$ against $21.3$ becomes $1.6$ against $36.0$.

% ------------------------------------- merged: methods in open-set + metric gameability
\subsection{Adaptation in the open-set setting, and a metric that rewards damage}
\label{sec:os-methods}

Table~\ref{tab:osmethods} runs the four adaptation methods of Section~\ref{sec:res-mech}
on the open-set stream, with the rejection rule held fixed at the source-calibrated
$q=5\%$ threshold so that only the adaptation differs between rows.

% source: src/exp25_openset_tta.py -> results/exp25_openset_tta.csv, 3 seeds
\begin{table*}[t]
\centering
\caption{Open-set TTA under the Tier-1 constraints, source-calibrated rejection at
$q=5\%$. OS$^{*}$ is mean per-class accuracy over known activities, UCDR the fraction of
truly-unknown windows rejected, H their harmonic mean. The two degenerate rules bound the
trade-off.}
\label{tab:osmethods}
\small
\begin{tabular}{lcccc|cccc}
\toprule
 & \multicolumn{4}{c|}{$\text{S}_1\!\to\!\text{S}_2$} & \multicolumn{4}{c}{$\text{S}_2\!\to\!\text{S}_1$} \\
Method & OS$^{*}$ & UCDR & H & \macf & OS$^{*}$ & UCDR & H & \macf \\
\midrule
Source-only & $0.6639$ & $0.3065$ & $\mathbf{0.4084}$ & $0.5216$ & $0.6717$ & $0.1427$ & $0.2276$ & $0.5769$ \\
AdaBN       & $0.3818$ & $0.2405$ & $0.2865$ & $0.3218$ & $0.4058$ & $0.1175$ & $0.1694$ & $0.3589$ \\
TENT        & $0.3826$ & $0.2381$ & $0.2849$ & $0.3220$ & $0.4065$ & $0.1170$ & $0.1690$ & $0.3592$ \\
DIGA        & $0.5322$ & $0.2178$ & $0.1419$ & $0.3856$ & $0.6572$ & $0.0480$ & $0.0891$ & $0.5450$ \\
PETAL       & $0.2812$ & $\mathbf{0.4375}$ & $0.3339$ & $0.2794$ & $0.3643$ & $\mathbf{0.2223}$ & $\mathbf{0.2551}$ & $0.3387$ \\
PeTTA       & $0.3312$ & $0.3299$ & $0.3208$ & $0.2962$ & $0.3820$ & $0.1617$ & $0.2148$ & $0.3429$ \\
\midrule
Reject none & $0.8494$ & $0.0000$ & $0.0000$ & $0.5552$ & $0.7529$ & $0.0000$ & $0.0000$ & $0.6072$ \\
Reject all  & $0.0000$ & $1.0000$ & $0.0000$ & $0.0000$ & $0.0000$ & $1.0000$ & $0.0000$ & $0.0000$ \\
\bottomrule
\end{tabular}
\end{table*}

In the direction with no measurable shift to recover, no method improves on the frozen
source model, and the ordering is the one Section~\ref{sec:res-mech} predicts from how
much source normalisation each method retains. Closing the label space does not change
that conclusion; it reproduces it on a harder task.

The reverse direction is where the metric becomes the finding. PETAL records the best
H-score of any method, $0.2551$ against source-only's $0.2276$, and it does so with
OS$^{*}$ of $0.3643$ against $0.6717$: it classifies the known activities roughly half as
well and still wins. It also records the best UCDR in both directions.

Neither is unknown detection. PETAL's normalisation channel damages the model globally, as
Table~\ref{tab:bnchannel} shows, which lowers its confidence on every window including the
known ones. More windows fall below a fixed threshold, UCDR rises mechanically, and
because H is a harmonic mean it rewards lifting whichever component was the bottleneck.
Unknown detection was the bottleneck, so a method that degrades uniformly is scored as a
method that detects novelty.

This is worth stating as a methodological result rather than an aside. UCDR has an obvious
trivial maximiser, which is why it is conventionally paired with OS$^{*}$; the pairing is
assumed to remove the problem. It does not. H, the standard open-set score and the one
LEAD and its comparators report, is itself maximised in part by uniform degradation
whenever the unknown component is the weaker of the two. An open-set result should
therefore be read with the known-class number beside it, and a method that wins on H while
losing substantially on OS$^{*}$ should be treated as a null rather than a positive. Every
H-score we report below is accompanied by its two components for this reason.

% ---------------------------------------------------------------- detector AUROC
\subsection{Confidence-based rejection is at or below chance}
\label{sec:os-detect}

Every method in Table~\ref{tab:osmethods} rejects on maximum softmax, which is the standard
choice. Whether that score carries any information about novelty on this data can be
asked directly, and independently of any threshold, by measuring how well it separates
known from unknown windows in the target stream.

% source: src/exp26_unknown_scores.py, src/exp28_auroc_ci.py -> results/exp28_auroc_ci.csv
\begin{table}[t]
\centering
\caption{AUROC for separating unknown from known target windows, averaged over 3 source
seeds, with $95\%$ bootstrap intervals over $2000$ resamples of the target stream.
$0.5$ is chance. All scores are computed from the frozen source model; none uses target
labels.}
\label{tab:osauroc}
\small
\begin{tabular}{@{}lcc@{}}
\toprule
Score & $\text{S}_1\!\to\!\text{S}_2$ & $\text{S}_2\!\to\!\text{S}_1$ \\
\midrule
Mahalanobis to prototype & $\mathbf{0.6377}$ $[0.620,0.655]$ & $\mathbf{0.5373}$ $[0.522,0.553]$ \\
Cosine to prototype      & $0.5668$ $[0.550,0.584]$ & $0.5070$ $[0.492,0.523]$ \\
\midrule
Maximum softmax          & $0.4952$ $[0.478,0.513]$ & $0.4434$ $[0.427,0.460]$ \\
Entropy                  & $0.4655$ $[0.447,0.483]$ & $0.4386$ $[0.422,0.455]$ \\
Maximum logit            & $0.4580$ $[0.441,0.475]$ & $0.4343$ $[0.418,0.451]$ \\
Free energy              & $0.4261$ $[0.409,0.443]$ & $0.4272$ $[0.411,0.443]$ \\
\bottomrule
\end{tabular}
\end{table}

The four confidence-derived scores are at or below chance. Entropy, maximum logit and free
energy have intervals lying entirely below $0.5$ in both directions. Maximum softmax, the
score the adaptation methods use, spans chance entering $\text{S}_2$
($[0.478,0.513]$) and lies below it entering $\text{S}_1$. None of the four is
usefully informative, and three are reliably anti-informative, meaning the model is on
average \emph{more} confident on the activities it was never trained on than on the ones it
was.

The direction of the effect is what makes it more than a null. A below-chance detector is
not a weak detector; it is a detector pointing the wrong way, and a system that rejects
low-confidence windows will preferentially discard the known activities it should be
keeping. The two prototype-distance scores, computed in the model's representation rather
than from its output, are the only ones above chance.

This parallels the per-sensor feature result of \cite{kbaier_closedset}, where features
the standard pipeline treats as its primary signal scored below stratified-random guessing
across buildings. In both cases a quantity the field treats as informative is not merely
uninformative on this data but carries signal of the wrong sign, and in both cases that is
only visible because a chance baseline was computed.

% ------------------------------------------- Mahalanobis, matched rate first
\subsection{Prototype distance helps, by about half what it first appears to}
\label{sec:os-maha}

Swapping maximum softmax for the Mahalanobis distance of Table~\ref{tab:osauroc} improves
unknown detection by a factor of $1.58$ over chance entering $\text{S}_2$ and $1.25$
entering $\text{S}_1$, measured at a rejection rate held equal between the two scores.
That is the number we would ask a reader to carry, and it is substantially smaller than
the one the natural comparison produces. We give both, because the gap between them is a
general hazard in open-set reporting rather than an accident of this experiment.

Compared at the same source-calibrated threshold $q=5\%$, the swap looks much larger: H
rises from $0.4084$ to $0.5176$ entering $\text{S}_2$ and from $0.2276$ to $0.4289$
entering $\text{S}_1$, gains of $10.9$ and $20.1$ points. Taken at face value this is the
largest effect anywhere in this study.

It is mostly an artefact of how much each score rejects. The two scores have different
distributions, so a threshold placed at the same percentile of \emph{source} validation
cuts the \emph{target} stream in very different places. Entering $\text{S}_2$, maximum
softmax rejects $31.9\%$ of the stream and the Mahalanobis distance rejects $54.0\%$,
against a true unknown share of $40\%$. The higher UCDR is bought by discarding half the
data, and OS$^{*}$ falls from $0.6639$ to $0.4447$ to pay for it. This is the same
mechanism as PETAL's apparent win in Section~\ref{sec:os-methods}, reached by a different
route, and it is why that section argues for reading H beside its components.

The fair comparison fixes the rejection rate and asks which score catches more true
unknowns. Table~\ref{tab:osmatched} sets every threshold so that exactly $20\%$ of the
target stream is rejected. Doing so requires a target quantile, so these are diagnostic
figures and not Tier-1 results; they are reported to size the detector, not to be deployed.

% source: src/exp27_proto_depth.py -> results/exp27_proto_depth.csv (matched20 rows)
\begin{table}[t]
\centering
\caption{UCDR with every threshold set to reject exactly $20\%$ of the target stream.
Chance is $0.2000$. DIAGNOSTIC: the matched rate uses a target quantile and is not
available without target data.}
\label{tab:osmatched}
\small
\begin{tabular}{@{}lcccc@{}}
\toprule
 & \multicolumn{2}{c}{$\text{S}_1\!\to\!\text{S}_2$} & \multicolumn{2}{c}{$\text{S}_2\!\to\!\text{S}_1$} \\
Score & UCDR & $\times$ chance & UCDR & $\times$ chance \\
\midrule
Mahalanobis, $64$-d embedding & $\mathbf{0.3162}$ & $1.58$ & $0.2342$ & $1.17$ \\
Mahalanobis, raw $13$-d input & $0.3125$ & $1.56$ & $0.2435$ & $1.22$ \\
Mahalanobis, $128$-d          & $0.2988$ & $1.49$ & $0.2031$ & $1.02$ \\
Mahalanobis, $256$-d          & $0.2451$ & $1.23$ & $\mathbf{0.2508}$ & $1.25$ \\
Maximum softmax               & $0.1980$ & $0.99$ & $0.1975$ & $0.99$ \\
\midrule
Chance                        & $0.2000$ & $1.00$ & $0.2000$ & $1.00$ \\
\bottomrule
\end{tabular}
\end{table}

The gain is real and it is modest. Prototype distance catches between $1.17$ and $1.58$
times as many unknowns as rejecting at random, and maximum softmax catches exactly as many
as random, to two decimal places, in both directions. The deployable configuration of
Section~\ref{sec:os-deploy} uses the $64$-d prototype distance at $q=5\%$ and obtains H
$0.5176$ and $0.4289$; those are the numbers a system would see, and the $1.58$ and $1.25$
figures are what the detector is actually contributing inside them.

We report this sequence in full rather than only its conclusion because the first
comparison we ran was the $q=5\%$ one, and it is the comparison an open-set paper would
normally present. Nothing in the protocol flags it as wrong. It took measuring the
rejection rate, which neither H nor UCDR exposes, to see that two scores held at the same
nominal threshold were operating $22$ percentage points apart.

% ---------------------------------------------------------------- threshold transfer
\subsection{A source-calibrated threshold does not transfer}
\label{sec:os-threshold}

The Tier-1 constraints require the rejection threshold to be fixed from source data, since
calibrating it on the target stream would use the labels the setting assumes are absent.
Section~\ref{sec:os-maha} showed that two scores held at the same nominal percentile
rejected very different fractions of the target stream. Table~\ref{tab:osthresh} shows
that this is systematic across the whole grid.

% source: src/exp27_proto_depth.py -> results/exp27_proto_depth.csv
\begin{table*}[t]
\centering
\caption{Fraction of the \emph{target} stream rejected by a threshold set to reject $q\%$
of the \emph{source} validation split. True unknown share is $40\%$ entering
$\text{S}_2$ and $30\%$ entering $\text{S}_1$.}
\label{tab:osthresh}
\scriptsize
\begin{tabular}{@{}lccccc|ccccc@{}}
\toprule
 & \multicolumn{5}{c|}{$\text{S}_1\!\to\!\text{S}_2$} & \multicolumn{5}{c}{$\text{S}_2\!\to\!\text{S}_1$} \\
Score & $q{=}1$ & $2.5$ & $5$ & $10$ & $20$ & $q{=}1$ & $2.5$ & $5$ & $10$ & $20$ \\
\midrule
Maximum softmax & $22.9$ & $26.5$ & $31.9$ & $37.7$ & $50.5$ & $7.3$ & $10.5$ & $14.8$ & $21.7$ & $36.5$ \\
Mahalanobis $64$-d  & $43.6$ & $49.0$ & $54.0$ & $60.0$ & $68.6$ & $22.1$ & $25.8$ & $31.0$ & $34.9$ & $48.9$ \\
Mahalanobis $128$-d & $62.8$ & $66.5$ & $71.2$ & $75.6$ & $80.5$ & $39.6$ & $47.3$ & $54.1$ & $60.5$ & $68.8$ \\
Mahalanobis $256$-d & $41.4$ & $54.2$ & $61.8$ & $72.2$ & $86.9$ & $49.6$ & $52.7$ & $55.9$ & $62.4$ & $67.4$ \\
Mahalanobis raw $13$-d & $54.2$ & $58.4$ & $69.5$ & $76.1$ & $87.8$ & $34.1$ & $37.7$ & $42.6$ & $56.2$ & $64.7$ \\
\bottomrule
\end{tabular}
\end{table*}

A threshold placed to reject one per cent of source validation rejects between $7.3\%$ and
$62.8\%$ of the target stream, the range depending only on which score it is placed on. At
$q=5\%$ the spread is $14.8\%$ to $71.2\%$. The nominal percentile, which is the quantity
an open-set method reports and tunes, carries almost no information about how aggressively
the system will actually reject once it is installed in another building.

The cause is the one running through this study. A threshold is a statistic of the source
distribution, and source statistics do not transfer across buildings. This is the property
responsible for two separate findings in the closed-set study \cite{kbaier_closedset}:
that replacing batch-norm statistics with the target stream's is harmful, and that
per-sensor features are anti-informative across buildings. It arrives here at a third
place, the rejection rule. There is nothing specific to our detector in it: any open-set or
universal method that fixes a threshold on source data inherits the problem, and any that
fixes one on target data has left the test-time setting.

We note the practical consequence without being able to solve it. A deployment cannot
observe its own rejection rate against ground truth, but it \emph{can} observe the rate
itself, which is label-free. A rejection rate far from the designer's intent is detectable
in the field even when correctness is not, and we suggest reporting it alongside H as a
matter of course. Calibrating a threshold that holds across domains, rather than merely
across resamples of one domain, we leave open.

% ---------------------------------------------------------------- representation depth
\subsection{Depth does not help, and the detector is leaving signal unused}
\label{sec:os-depth}

The prototypes of Section~\ref{sec:os-maha} are built in the model's final
$64$-dimensional embedding, which is trained to separate three known activities and
therefore to discard whatever does not serve that. Whether novelty survives the
compression is testable at no cost, by building the same prototypes at three other
depths, including none at all.

% source: src/exp27_proto_depth.py -> results/exp27_proto_depth.csv
\begin{table}[t]
\centering
\caption{AUROC for unknown-versus-known by the representation the prototype distance is
computed in, with the supervised reference of Table~\ref{tab:osceiling} for scale. Chance is
$0.5$.}
\label{tab:osdepth}
\small
\begin{tabular}{@{}lcc@{}}
\toprule
Representation & $\text{S}_1\!\to\!\text{S}_2$ & $\text{S}_2\!\to\!\text{S}_1$ \\
\midrule
Raw $13$-d input, no network & $\mathbf{0.6392}$ & $0.5253$ \\
Final embedding, $64$-d      & $0.6377$ & $\mathbf{0.5373}$ \\
Second block, $128$-d        & $0.6198$ & $0.5076$ \\
First block, $256$-d         & $0.5982$ & $0.5107$ \\
\midrule
Maximum softmax              & $0.4952$ & $0.4434$ \\
Supervised reference (diagnostic)& $0.7354$ & $0.6327$ \\
\bottomrule
\end{tabular}
\end{table}

The raw input features match the final embedding, $0.6392$ against $0.6377$ in one
direction and $0.5253$ against $0.5373$ in the other, differences well inside the seed
spread, and the intermediate blocks are worse than both. Passing the window through a
trained network neither adds novelty information nor destroys it, so the rejection
component of Section~\ref{sec:os-deploy} does not require the classifier at all and could
run beside a model of any architecture.

It is tempting to read that as a limit of the features: if a representation trained on the
known classes carries no more novelty signal than the untrained input, perhaps thirteen
window statistics simply do not encode what separates an unseen activity. We tested that
reading rather than asserting it, and it is wrong.

% source: src/exp29_feature_limit.py -> results/exp29_feature_limit.csv
\begin{table}[t]
\centering
\caption{Is novelty present in the $13$-d description? Per-feature AUROC separating
unknown from known windows, and a supervised known-versus-unknown classifier trained on
the target building under GroupKFold over recording day. The classifier uses target
labels and is a DIAGNOSTIC upper bound, not a method.}
\label{tab:osceiling}
\small
\begin{tabular}{@{}lcc@{}}
\toprule
 & $\text{S}_2$ & $\text{S}_1$ \\
\midrule
Features differing after Benjamini-Hochberg at $5\%$ & $13/13$ & $13/13$ \\
Strongest single feature & $0.6972$ & $0.3700$ \\
\quad which one & \texttt{t\_time\_cos} & \texttt{g\_max\_active} \\
Supervised reference, logistic & $0.6620$ & $0.5701$ \\
Supervised reference, gradient boosting & $\mathbf{0.7354}$ & $\mathbf{0.6327}$ \\
\midrule
Our label-free detector (Table~\ref{tab:osdepth}) & $0.6377$ & $0.5373$ \\
\bottomrule
\end{tabular}
\end{table}

All thirteen marginals differ between the withheld and retained activities, in both
buildings, after correcting for the thirteen tests. A supervised classifier given the same
features reaches AUROC $0.7354$ and $0.6327$. The information is present, and our
label-free detector reaches about ten AUROC points below what a labelled one extracts from
the identical description.

The single most informative feature is worth noting. In $\text{S}_2$ it is
\texttt{t\_time\_cos}, the cosine of time of day, at $0.6972$ on its own, which is higher
than our entire multivariate detector. The withheld activities are toileting and cleaning,
and they occur at different hours than meal preparation, television and music. A
Mahalanobis distance to class prototypes treats all thirteen dimensions as one geometry
and dilutes that; it does not exploit the one axis that most separates the groups.

So the honest statement has two parts rather than one. The learned representation adds
nothing over the raw features, which is the clean negative result the depth sweep
establishes. But the features are not the binding constraint: a labelled detector on the
same thirteen numbers does roughly ten points better, so there is headroom a better
label-free score could reach, and the most obvious place to look for it is temporal
structure rather than a distance in feature space. Section~\ref{sec:os-temporal} follows
that up. The labelled figures should be read as a reference point rather than as a ceiling. They
are what one particular supervised learner extracts, and
Section~\ref{sec:os-temporal} reaches $0.6616$ entering $\text{S}_1$ with an unsupervised
distance on four features, against the $0.6327$ obtained there with gradient boosting on
all thirteen and target labels. A bound that an unsupervised method beats was not a bound;
it bounds that learner on that feature set, and the true limit of this description of the
data is higher than either number. We keep the comparison because the direction of the gap
is informative --- a labelled detector does better than ours on the identical features, so
signal is being left unused --- while noting that its size is not a reliable estimate of
how much.

% ---------------------------------------------------------------- temporal detector
\subsection{A better detector does not produce a better system}
\label{sec:os-temporal}

Section~\ref{sec:os-depth} located roughly ten AUROC points of unused signal and pointed
at time of day. Acting on that raises a question of method before one of performance: we
identified the temporal axis by looking at target labels, so a detector built on it is
selected with target information through the analyst rather than through the code. Scoring
it as Tier-1 would be the error we document in others, arrived at by a different route.

The test is whether a rule that sees only source data would have chosen the same axis.
Hold out each known activity from the source prototypes in turn, score that activity's own
\emph{source} windows as though they were novel, and record which dimensions detect them;
averaging over the three known classes gives per-dimension weights that use no target data
at all. The rule selects the temporal axis in both directions. \texttt{t\_time\_sin} and
\texttt{t\_time\_cos} take the largest weights, and the four temporal dimensions average
$1.58$ and $1.96$ against $0.74$ and $0.57$ for the nine activity statistics. The axis is
therefore selectable without target labels on this data, which is what distinguishes a
usable rule from an observation we happened to make with the target labels in front of us.

The scope of that claim needs stating plainly, because it is the strongest interpretive
step in this paper. What we have shown is that one source-only rule selects the same axis
that target labels identify, in two transfer directions between two apartments of one
dataset, for one choice of withheld activities. Whether leave-one-class-out weighting
recovers the discriminative axis in general --- in other buildings, other sensor layouts,
other label splits, or where novelty is not separated by time of day --- is untested here
and does not follow from these two agreements. The rule could equally be selecting
temporal features because cyclical time is the one part of this feature vector not derived
from sensor activations, in which case it would transfer, or because toileting and
cleaning happen at distinctive hours in these two particular homes, in which case it would
not. We cannot distinguish those with two buildings, and a reader should treat the design
rule as a hypothesis with source-only support rather than as an established method.

% source: src/exp30_temporal_detector.py -> results/exp30_temporal_detector.csv
\begin{table*}[t]
\centering
\caption{Detectors that isolate or upweight the temporal axis. $\Delta$ is a paired
bootstrap over the target stream against the $13$-d mean baseline, $2000$ resamples.
Variants marked [ph] select the axis using target
knowledge: they are DIAGNOSTIC upper bounds in the sense of
Section~\ref{sec:os-setup}, not deployable results. [t1] variants use source data only. H and the rejection rate are at the fixed $q=5\%$
source-calibrated threshold.}
\label{tab:ostemporal}
\small
\begin{tabular}{lccc|ccc}
\toprule
 & \multicolumn{3}{c|}{$\text{S}_1\!\to\!\text{S}_2$} & \multicolumn{3}{c}{$\text{S}_2\!\to\!\text{S}_1$} \\
Detector & AUROC & $\Delta$ & H / rej & AUROC & $\Delta$ & H / rej \\
\midrule
$13$-d mean [t1]        & $0.6392$ & --- & $0.4957$ / $70\%$ & $0.5253$ & --- & $0.4279$ / $43\%$ \\
$13$-d max-dim [t1]     & $0.6308$ & $-0.009$ & $\mathbf{0.5235}$ / $57\%$ & $0.5745$ & $+0.049$ & $0.4410$ / $54\%$ \\
LOCO-weighted [t1]      & $0.6373$ & ns & $0.4069$ / $74\%$ & $0.5888$ & $+0.064$ & $\mathbf{0.4825}$ / $55\%$ \\
Temporal $4$-d [ph]     & $0.6389$ & ns & $0.4685$ / $23\%$ & $\mathbf{0.6616}$ & $+0.136$ & $0.4986$ / $25\%$ \\
Time-of-day $2$-d [ph]  & $0.5915$ & $-0.047$ & $0.5053$ / $26\%$ & $0.6266$ & $+0.101$ & $0.4207$ / $16\%$ \\
\bottomrule
\end{tabular}
\end{table*}

As a detector the temporal axis works, and it works only where there is a shift to detect.
Entering $\text{S}_1$, every variant separates from the baseline: the source-only
weighting gains $0.064$ AUROC and the post-hoc temporal distance $0.136$. Entering
$\text{S}_2$, where the cross-building penalty is $1.6$ points, nothing improves and two
variants are significantly worse. This is the same asymmetry the adaptation methods
followed in Section~\ref{sec:os-methods}, now in a component that does not adapt anything.

As a system it changes nothing, and that is the more useful result. At the fixed
source-calibrated threshold the source-only weighting rejects $74\%$ of the forward stream
and its H falls to $0.4069$, below the $0.5176$ of the configuration in
Section~\ref{sec:os-deploy}. The better-separating score is the worse-calibrated one. The
$13$-d max-dim variant shows the converse: it records the best H in both directions while
its AUROC is \emph{significantly worse} in one of them, so its apparent advantage is a
rejection-rate effect of the kind Section~\ref{sec:os-maha} dissects rather than better
detection. We decline to recommend it for that reason.

What this establishes is a constraint rather than a method. Detection quality and
deployed open-set performance come apart here: a score can gain $0.064$ AUROC and lose
$0.11$ H at the threshold a deployment would actually use. Until a threshold can be
calibrated so that it transfers, improving the detector does not improve the system, and
the limiting factor in Table~\ref{tab:osthresh} is not the one anybody is working on.

% ---------------------------------------------------------------------------
%  Formal table of where the reference open-set implementation leaves the
%  test-time setting. Same format as the paper/code divergence table.
%  Source: github.com/ispc-lab/LEAD, inspected at HEAD.
% ---------------------------------------------------------------------------
\subsection{Where the reference implementation leaves the test-time setting}
\label{sec:os-lead}

The open-set and universal domain adaptation methods we build on are distributed with
implementations that assume offline access to the target domain. That assumption is not
stated as a limitation, and it is not visible from the papers, but it decides what the
reported numbers mean. Table~\ref{tab:leadviol} records the four places where the
reference implementation of LEAD \cite{lead} departs from the Tier-1 definition of
Section~\ref{sec:method-tiers}, with the line that does it.

\begin{table*}[t]
\centering
\caption{The reference LEAD implementation (\texttt{ispc-lab/LEAD}) against the Tier-1
test-time constraints. Line numbers are from \texttt{train\_target.py} unless marked.
Together these make the published protocol source-free domain adaptation rather than
test-time adaptation: it requires the whole target set in advance, revisits it fifty
times, and evaluates on the data it adapted on.}
\label{tab:leadviol}
\small
\begin{tabular}{@{}p{0.20\textwidth}p{0.17\textwidth}p{0.26\textwidth}p{0.29\textwidth}@{}}
\toprule
Mechanism & Code & Constraint violated & Consequence \\
\midrule
Multi-epoch offline training on the target set
 & \texttt{epochs=50} (\texttt{config/model\_config.py}:15); loop at :375
 & Adaptation must be a single online pass over the stream
 & The whole target set must exist before adaptation starts, and every sample is revisited
   fifty times. No deployed system has this. \\
\addlinespace[3pt]
Adaptation and evaluation share one dataset
 & :360--362, both \texttt{DataLoader}s wrap the same \texttt{target\_dataset}
 & Evaluation must not be adapted on
 & Every sample scored at :383 was trained on at :377. The reported score is training
   performance, and no held-out target split exists anywhere in the file. \\
\addlinespace[3pt]
Pseudo-labels built from a full pass over the evaluation set
 & :209 calls \texttt{obtain\_LEAD\_pseudo\_labels} on \texttt{test\_dataloader} at the
   start of every epoch
 & Predictions must depend only on samples already seen
 & The model sees the entire evaluation set, including its own future, before adapting.
   Known- and unknown-space bases are fitted on data not yet encountered. \\
\addlinespace[3pt]
Stream order discarded
 & \texttt{shuffle=True} on the adaptation loader, :360
 & The target stream must be processed in acquisition order
 & Temporal structure is destroyed and each batch mixes the whole recording period, so
   batch statistics are those of the full target distribution rather than of the moment. \\
\bottomrule
\end{tabular}
\end{table*}

None of the four is a defect of the method as described; each is a property of the
evaluation harness it ships with. The distinction matters because the first and third in
particular make the reported figures an upper bound that a streaming deployment cannot
approach, and because a reimplementation that honoured the constraints would not be
comparable with the published numbers. All open-set results in
Section~\ref{sec:os-methods} are therefore reported against our own Tier-1 source-only
baseline rather than against published figures obtained under this harness.


% ============================================================================
\section{Discussion}
% ---------------------------------------------------------------- discussion
\subsection{Discussion}
\label{sec:os-discussion}

Adding a label-space shift to a geographic one does not give adaptation more to work with.
The cross-building penalty is $1.6$ points entering $\text{S}_2$ and $36.0$ entering
$\text{S}_1$ once classes are withheld, a wider asymmetry than the closed-set
$0.6$ and $21.3$, and in neither direction does any of the four methods improve on a frozen
source model once H is read beside its components. The combined shift is harder, not
richer.

What does help is replacing the rejection score, and only modestly: between $1.17$ and
$1.58$ times chance at a matched rejection rate, against a maximum softmax baseline that
sits exactly at chance. This is a consistent reading alongside
Section~\ref{sec:res-mech}, where DIGA's prototype classifier was the one component that
helped without touching normalisation statistics, but we did not isolate that component in
the open-set setting and the two results are not a replication of one another. What they
have in common is the direction of the leverage: in both, the useful part was a distance
in representation space, and the parts that manipulate source statistics were harmful or
inert.

Three things we would want a reader to take from this section are not about our system.

Metrics first. UCDR has a trivial maximiser and is conventionally paired with OS$^{*}$ to
remove it. That pairing is insufficient, because H is itself partly maximised by uniform
degradation whenever unknown detection is the weaker component, which it usually is.
PETAL records the best H in one direction while classifying the known activities half as
well as doing nothing. Open-set results should be reported with both components visible.

Thresholds second. A threshold set to reject $1\%$ of source validation rejects between
$7\%$ and $63\%$ of a target stream, depending only on which score carries it. Any
open-set method that calibrates on source inherits this, and any that calibrates on target
has left the test-time setting. We think the rejection rate, which is observable without
labels, belongs in open-set reporting next to H. How to calibrate a threshold that holds
across domains rather than across resamples of one is open, and is the direction we would
pursue next.

Chance baselines third. The two results in this section that are most likely to generalise
are both negative, and neither is visible without a chance baseline: four standard
confidence scores are at or below chance as novelty detectors, and a learned representation
carries no more novelty signal than the untrained input. Neither would appear as a problem
in a table of H-scores alone.

The limits are the same as the closed-set study's. Two buildings, one dataset, one choice
of withheld activities. The asymmetry we report is a property of this pair, and the
specific values of $1.6$ and $36.0$ belong to it. What we would expect to transfer is the
procedure: measure the penalty per cell against a target-trained model at matched sample
size before running any method, report rejection rates beside H, and compute a chance
baseline for any novelty score before believing it.


\section{Conclusion}
\label{sec:os-conclusion}

Adding a label-space shift to a geographic one makes cross-building activity recognition
harder without giving adaptation more to work with. Across a $2\times2$ grid of transfer
direction and label space, the methods track the measured penalty rather than the task:
they help only where there is a real shift, and in the direction where the penalty is
$1.6$ points none of them improves on doing nothing.

What survives is small and does not adapt. A frozen source model with a Mahalanobis
rejection rule, with prototypes and threshold fixed on source data before the target
stream is read, reaches H-scores of $0.5176$ and $0.4289$; every component we removed to
reach it --- the four adaptation methods, the confidence-based rejection scores, and the
learned embedding, which a distance on the raw thirteen features matches --- was worth
removing.

Three findings seem to us likely to outlive this dataset, and none of them is about
our system.

The standard open-set score can be raised by degrading a model uniformly whenever
unknown detection is the weaker of its two terms, as it was in both directions here, so
it should never be reported without its known-class component.

Four standard confidence-based novelty scores are at or below chance here, which was
visible only because a chance baseline was computed, and is a reason to compute one
before trusting any detector in a new domain.

And a threshold calibrated on source data did not transfer in either direction here: the
same nominal percentile rejects $7\%$ or $63\%$ of a target stream depending on the
score it sits on, and improving the detector underneath did not help while that was
true. We expect the mechanism to be general, since a threshold is a statistic of the
source distribution, but we have measured it on two buildings.

The limits are those of the study it extends. Two apartments of one dataset, one choice of
withheld activities, and a temporal axis that turns out to carry much of the novelty
signal in these particular homes. The specific numbers belong to that pair. What we would
expect to transfer is the procedure: measure the penalty per cell against a target-trained
model at matched sample size before running any method, compute a chance baseline for any
novelty score, report the rejection rate beside the H-score, and treat a method that wins
on H while losing on known-class accuracy as a null. Calibrating a rejection threshold
that holds across domains rather than across resamples of one is, on this evidence, the
problem worth working on next.


% ============================================================================
\section*{Reproducibility}

Code and every result file are available at
\url{https://github.com/mariemkbaier/TTA-Test-Time-Adaptation}. Each table above names
the script and result file that produced it, and
\texttt{src/verify\_numbers.py} checks every numeric claim in this manuscript against
those files.

\section*{Acknowledgements}

This research was supported by the Natural Sciences and Engineering Research Council of
Canada (NSERC).

% ============================================================================
\begin{thebibliography}{99}

\bibitem{kbaier_closedset} M. Kbaier, O. Abderrahim, M. Amayri, N. Bouguila, Sensor
identity does not transfer: revisiting test-time adaptation for cross-building human
activity recognition, under review.

\bibitem{lead} S. Qu, T. Zou, L. He, F. R\"ohrbein, A. Knoll, G. Chen, C. Jiang, LEAD:
learning decomposition for source-free universal domain adaptation, in: Proceedings of
the IEEE/CVF Conference on Computer Vision and Pattern Recognition, 2024,
pp. 23334--23343.

\bibitem{tapia2004} E. M. Tapia, S. S. Intille, K. Larson, Activity recognition in the
home using simple and ubiquitous sensors, in: Pervasive Computing, Springer, Berlin,
Heidelberg, 2004, pp. 158--175.

\end{thebibliography}

\end{document}
